% !TeX root = main.tex
\chapter{特征值与特征向量}
特征值与特征向量，以及接下来的广义特征分解，都是为了分解算子，本质上是在研究线性空间上自同态的结构。

\section{特征值}
% TODO: https://www.zhihu.com/question/37615080/answer/1982559561807307059
\begin{definition}[算子]
    算子：从一个线性空间到自身的映射
\end{definition}

算子在定义上已经具备了可逆的基础。如果线性映射是不同线性空间之间的同态，那算子就是自同态。
算子比一般的线性映射有趣的地方在于： 算子可以求幂。

预想研究算子的性质，一个很好的思路是把算子分解成一个个不变子空间上算子的组合。
特征向量就是最简单的不变子空间——\textbf{一维不变子空间}里的元素。在这个子空间上，算子被简化为纯粹的拉伸与压缩。
\[
    Tv=\lambda v \Leftrightarrow \nullspace(T-\lambda I)\neq {0}
\]

不同特征值的分解下产生的不变子空间相互``正交''\footnote{这里的正交仅指子空间构成直和}。
以下比较\textit{Linear Algebra Done Right} 与\textit{高等代数}中证明方法的不同：

\begin{theorem}[不同特征值对应特征向量无关]
    设 \(T \in \L(V)\)那么分别对应于\(T\) 的不同特征值的特征向量构成的每个组都线性无关
\end{theorem}

\begin{proof}
    （\cite[p.~114]{LinearAlgebraDone2024}）
    假设欲证结论不成立。那么存在最小的正整数 \(m\)， 使得\(T\)
    对应于其互异特征值\(\lambda_1, \dots ,\lambda_{m}\)
    的特征向量 \(v_1, \dots, v_m\)构成线性相关向量组。
    于是存在\(\alpha_1, \dots ,\alpha_{m}\)，其中各数均非零，使得
    \[\alpha_1 v_1 + \cdots + \alpha_m v_m = 0\]
    将 \(T-\lambda_{m}I\) 作用于上式两侧，得
    \[
        \alpha_1 (\lambda_1 - \lambda_m)v_1 + \cdots +
        \alpha_{m-1} (\lambda_{m-1} - \lambda_m)v_{m-1} = 0
    \]
    因为特征值\(\lambda_1,\dots ,\lambda_{m}\)互异，所以上式中各系数均不等于0。
    于是 \(v_1, \dots , v_{m-1}\) 就是由对应于 \(T\)
    的互异特征值的 \(m- 1\) 个特征向量构成的线性相关向量组，这和 \(m\) 最小的假设相矛盾． 此矛盾
    就说明欲证命题成立．
\end{proof}

\begin{proof}
    （\cite[p.~204 定理~8]{GaoDengDaiShu2019}）
    对特征值的个数做数学归纳法。由于特征向量是不为零的。单个特征向量显然线性无关。先设属于\(k\)
    个不同特征值的特征向量线性无关，我们证明属于\(k+1\) 个不同特征值\(\lambda_1,\dots
    ,\lambda_{k+1}\) 的特征向量\(\xi_1,\dots, \xi_{k+1}\) 也线性无关。
    设关系式
    \begin{equation}
        a_1\xi_1 + a_2\xi_2 + \dots + a_{k+1}\xi_{k+1} = 0 \label{1}
    \end{equation}
    两边同乘\(\lambda_{k+1}\) 有：
    \begin{equation}
        a_1\lambda_{k+1}\xi_1 + a_2\lambda_{k+1}\xi_2 +
        \cdots + a_{k+1}\lambda_{k+1}\xi_{k+1} = 0 \label{eq:2}
    \end{equation}
    式子两端同时施加变换\(\A\)，有：
    \begin{equation}
        a_1\lambda_1\xi_1 + a_2\lambda_1\xi_2 + \cdots +
        a_{k+1}\lambda_1\xi_{k+1} = 0 \label{eq:3}
    \end{equation}
    \cref{eq:2} 减去 \cref{eq:3}，得：
    \[
        a_1(\lambda_1 - \lambda_{k+1})\xi_1 + a_2(\lambda_1
        - \lambda_{k+1})\xi_2 + \cdots +
        a_{k}(\lambda_{k} - \lambda_{k+1})\xi_{k+1} = 0
    \]
    根据归纳假设，\(\xi_1,\dots \xi_{k}\) 线性无关，于是
    \[
        a_{i}(\lambda_{i} - \lambda_{k+1}) = 0, \quad i=1,\dots ,k
    \]
    但\(\lambda_{i}-\lambda_{k+1}\neq 0\)，所以\(a_{i}=0\)。
    从而\(a_{k+1}\) 也为零。这说明\(\lambda_1,\dots ,\lambda_{k+1}\) 线性无关
\end{proof}

短短的几行证明，不知蕴含了作者大一多少不解\UseVerb{cry}。哪个证明更加符合直觉，
想必是高下立判了吧\UseVerb{surprised}\footnote{其实回头看来，二者区别也没有那么大：无非变换顺序的问题。
但初学时确实会产生不小困惑。}

不同特征值对应特征向量线性无关这一定理，还可以用于证明向量线性无关：
\begin{example}
    考虑\(V=\Span(\e^{\lambda_{1}x},\dots \e^{\lambda_{n}x}  )\)。
    不难发现\(\e^{\lambda_{i}x} \)是 \(Df = f'\) 算子的特征向量，
    对应特征值为 \(\lambda_{i}\)。
    因此\(\left\{ \e^{\lambda_{1}x},\dots \e^{\lambda_{n}x} \right\}\) 线性无关。
\end{example}

\(T\) 的不变子空间包括：
\begin{itemize}
    \item \(\nullspace T\) 和\(\range T\)，
        包括所有\(U\subseteq \nullspace T\) 和 \(V \supseteq \range T\)
    \item \hl{\(\nullspace p(T)\) 与 \(\range p(T)\)}
\end{itemize}
% todo: 补充可交换的多项式

特征值是算子本身的一种特质：
\begin{table}[htbp]
    \centering
    \begin{tabular}{cccccccccc}
        \hline
        & \(T\) & \(p(T)\) & \(T^{-1}\) & \(T^{\mathsf{T}}\) & \(P^{-1}TP\) & \(T_{\C}\) & \(T/U\) & \(T^*\) & \(aE - A\)\\
        \hline
        特征值 & \(\lambda\) & \(p(\lambda)\) & \(\lambda^{-1}\) & \(\lambda\) & \(\lambda\) & \(\lambda\) & \(\lambda\) & \(\tfrac{\abs{A}}{\lambda}\) & \(a-\lambda\)\\
        特征向量 & \(\xi\) & \(\xi\) & \(\xi\) & & \(P^{-1}\xi\) & & & \(\xi\) & \(\xi\)\\
        \hline
    \end{tabular}
\end{table}
\(T/U\) 表示商算子，\((T/U)v=Tv+U\)。

整个线性代数的计算可以说都建立在解方程上，求特征值也不例外。特征值的概念和矩阵零空间息息相关。

欲求\(T\xi = \lambda\xi\)，即解\((\lambda I - T)\xi = 0\)，
也就是求 \(\nullspace(\lambda I - T)\)。这个方程的意义就是试根：踏破铁鞋无觅处，
寻找满足存在 \((\lambda I - T)v\)的那个\(\lambda \)。
从这个角度上看，\(0\)是不可逆算子天然的特征值，
\(T\)的特征值个数不会超过 \(1+\dim \range T\)\footnote{更多计算可见\cref{theorem:matrix-determinent-lemma}}。

\section{最小多项式}

% thm:特征值的存在性（Tv一列必线性相关）
% thm:最小多项式的存在性（把整个的约化成小的，然后小的必然存在）
% 后一个定理的不平凡之处就在于，前者只证明了存在一个多项式对于某个特定向量v可以零化，后者则证明了整个算子都存在一个多项式可以零花
% (唯一性相当好证，略)
% todo:最小多项式计算方法
% 此时可能会有同学发问了？这和课本上写的不一样呐？不应该计算det lambda I - T，然后计算特征值吗？
% 至少在LADR这个框架下，应该是先有特征值和特征向量，然后有最小多项式。根据广义特征分解给出特征多项式，然后才发现嘿这不就是det lam I - A ?
% 先入为主的要求特征多项式det lam I - A去匹配广义特征分解就是抽象的，这时候和行列式还没啥关系呢。起码在此刻的叙事下，det的特征多项式更难

% thm 最小多项式的根是特征值and证明
% thm 最小多项式整除所有零化多项式（也非常好证，很经典）

% 商算子的最小多项式是原算子的最小多项式除以不变子空间的最小多项式？

% 为什么要研究最小多项式
% 这个东西也很复杂。简单来说你都可以从range(p(T))是T不变子空间入手说明，甚至在式子两边取行列式得到特征多项式。
% 但似乎这个研究与矩阵生成的代数等有很深刻联系，可见cherry studio笔记
\section{上三角矩阵}
算子的上三角矩阵看似很弱，但这一分解足够普遍也足够简单，为后文的对角化和广义特征分解打下了基础，还是得牢牢掌握。
% thm 上三角阵三个条件（挺直观的，不必证明）
% thm 5.40
% thm 5.41 上个定理只是说(T-lam I)...v可以被零化，这里直接说肯定有一维度必被零化
% 上三角阵充要条件！这个和后面很多证明一脉相承
\section{可对角化算子}
% 特征空间定义、三个性质（不必证）
% 可对角化充要条件

\begin{theorem}
    设 \(V\) 是有限维的，且 \(T \in \L(V)\). 那么 \(T\) 是可对角化的，
    当且仅当 \(T\) 的最小多项式等于 \((z - \lambda_1) \cdots (z - \lambda_m)\), 其中 \(\lambda_1, \ldots, \lambda_m \in \F\) 为互不相同的数。
\end{theorem}

\begin{proof}
    首先设 \(T\) 是可对角化的. 那么存在由 \(T\) 的特征向量构成的 \(V\) 的基 \(v_1, \ldots, v_n\). 令 \(\lambda_1, \ldots, \lambda_m\) 是 \(T\) 的所有互异特征值. 那么对每个 \(v_j\), 存在 \(\lambda_k\) 使得 \((T - \lambda_k I)v_j = 0\). 于是
    \[(T - \lambda_1 I) \cdots (T - \lambda_m I)v_j = 0,\]

    意味着 \(T\) 的最小多项式等于 \((z - \lambda_1) \cdots (z - \lambda_m)\).

    为了证明另一方向的蕴涵关系，
    现设 \(T\) 的最小多项式等于 \((z - \lambda_1) \cdots (z - \lambda_m)\), 其中 \(\lambda_1, \ldots, \lambda_m \in \mathbf{F}\) 为互不相同的数. 于是
    \[(T - \lambda_1 I) \cdots (T - \lambda_m I) = 0\]

    我们将通过对 \(m\) 用归纳法来证明 \(T\) 是可对角化的. 首先设 \(m = 1\). 那么 \(T - \lambda_1 I = 0\), 这意味着 \(T\) 是恒等算子的标量倍，
    也就说明了 \(T\) 是可对角化的.

    现在设 \(m > 1\) 且欲证结论在 \(m\) 为更小的值时都成立. \(\range(T - \lambda_m I)\) 这个子空间在 \(T\) 下是不变的。
    那么限制在 \(\range(T - \lambda_m I)\) 上的 \(T\) 就是 \(\range(T - \lambda_m I)\) 上的算子。

    若 \(u \in \range(T - \lambda_m I)\), 那么对某个 \(v \in V\) 有 \(u = (T - \lambda_m I)v\)，
    可得
    \[(T - \lambda_1 I) \cdots (T - \lambda_{m-1} I)u = (T - \lambda_1 I) \cdots (T - \lambda_m I)v = 0\]

    因此，
    \((z - \lambda_1) \cdots (z - \lambda_{m-1})\) 是限制于 \(\range(T - \lambda_m I)\) 上的 \(T\) 的最小多项式的多项式倍（由 5.29）。
    于是，由归纳假设，\(\range(T - \lambda_m I)\) 中存在由 \(T\) 的特征向量组成的基.

    设 \(u \in \range(T - \lambda_m I) \cap \nullspace(T - \lambda_m I)\). 那么 \(Tu = \lambda_m u\). 结合式 (5.64) 可得
    \begin{align}
        0 &= (T - \lambda_1 I) \cdots (T - \lambda_{m-1} I)u \\ &= (\lambda_m - \lambda_1) \cdots (\lambda_m - \lambda_{m-1})u.
    \end{align}

    因为 \(\lambda_1, \ldots, \lambda_m\) 是互异的，
    所以上式表明 \(u = 0\). 因此 \(\range(T - \lambda_m I) \cap \nullspace(T - \lambda_m I) = \{0\}\).

\end{proof}
